%% florilegium-book.tex -- the grown composition in a two-sided document: the garden is in
%% the outer margin of every page (left on even pages, mirrored on the right on odd ones).
%% LuaLaTeX, one pass:
%%   lualatex florilegium-book.tex
%% XeLaTeX works too (twice, the margins of the pages are known after the first run):
%%   xelatex --shell-escape florilegium-book.tex
%%   xelatex florilegium-book.tex && texlua florilegium-xe.lua florilegium-book && xelatex florilegium-book.tex
%% The option book needs a two-sided layout (twoside in the class and in geometry) and
%% works with the authored composition too: \usepackage[book]{florilegium}.
\documentclass[11pt,a4paper,twoside]{article}
\renewcommand\rmdefault{cmr}\renewcommand\sfdefault{cmss}\renewcommand\ttdefault{cmtt}
\usepackage[T1]{fontenc}
\usepackage{amsmath,amssymb,amsthm}
\usepackage[a4paper,twoside,inner=22mm,outer=68mm,top=32mm,bottom=30mm,headsep=8mm]{geometry}
\usepackage{fancyhdr}
\usepackage[grown,book]{florilegium}

\pagestyle{fancy}\fancyhf{}
\fancyhead[RO,LE]{\small\textsc{3\quad Procedural ornament}}
\fancyfoot[RO,LE]{\small\thepage}
\renewcommand\headrulewidth{0pt}

\theoremstyle{plain}
\newtheorem{proposition}{Proposition}[section]
\theoremstyle{definition}
\newtheorem{definition}[proposition]{Definition}
\newtheorem{construction}[proposition]{Construction}
\newtheorem{remark}[proposition]{Remark}
\newtheorem{exercise}[proposition]{Exercise}

\begin{document}
%% text-methods.tex -- shared body text: how the margin is drawn
\setcounter{section}{3}\setcounter{subsection}{1}
\subsection{The garden in the margin}

Everything in the left margin of this page is drawn while the page is being shipped out,
from a few hundred lines of Lua run by Lua\TeX\ (under Xe\LaTeX, by \texttt{texlua} beside it). There are no image
files, no imported paths and no external assets; every petal is a parametric curve evaluated for this page
and written straight into it as PDF path operators. This section records how it is done,
so that the margin documents itself.

\begin{definition}
A \emph{primitive} is a closed outline $\gamma\colon[0,1]\to\mathbb{R}^2$ authored by hand
together with a small set of free parameters $(p_1,\dots,p_k)$ and a recipe for painting
it. An \emph{instance} of the primitive samples the parameters from fixed intervals and
paints $\gamma$ in a local frame chosen by the caller. The composition on this page is
\emph{authored}; only the parameters are \emph{sampled}. This is what ``procedural'' means
here: no flower is grown from a grammar, but no two petals share a shape.
\end{definition}

\begin{construction}[Palette]
There is one document color, \texttt{ink}, and from it a ramp $t\mapsto\operatorname{tone}(t)$
from the darkest navy at $t=0$ to a pale mist at $t=1$. The ramp is interpolated in
OKLCH between six stops whose lightness runs from $0.25$ to $0.92$ and whose chroma peaks
at $0.175$ for the cobalt at $t=0.38$; the hue is that of \texttt{ink}, drifting $6^\circ$
toward violet in the darks and $8^\circ$ toward cyan in the lights, as ultramarine does when
it is thinned with water, and the chroma scales with that of \texttt{ink}. Mixing the
document color with white, as \texttt{ink!22} would, desaturates toward gray instead;
the perceptual ramp keeps the light end blue. A second color, \texttt{accent}, is a copper
used only for hairlines.
\end{construction}

\begin{construction}[Petal]
A petal of length $1$ pointing along $+y$ has a spine $x=b\,v^2$, $y=v$ ($b$ bends the
tip) and the obovate half-width
\[
  w(v)=\begin{cases}
    W\bigl(0.16+0.84\,\sin^{1.3}\frac{\pi v}{1.24}\bigr), & v<0.62,\\[2pt]
    W\sqrt{1-q^2},\quad q=\frac{v-0.62}{0.38}, & v\ge 0.62,
  \end{cases}
  \qquad W\in[0.40,0.48],
\]
a narrow claw, widest at $0.62$, with an elliptical cap. The margin is moved along the
normal by $0.03\,v\,n(6v)$, where $n$ is a three-octave gradient noise seeded per petal, so
it wavers more toward the tip; four petals in ten also get a shallow notch. The outline is
sampled once and used twice: as a clipping path, and as the stroke that gives the rim.
Inside the clip the petal is painted like a watercolor, in the order a painter would: a
flat wash at $\operatorname{tone}(t_0)$, $t_0\in[0.30,0.46]$; eight stacked white blobs at
the base, each at $10\%$ opacity, for the pale heart; fifteen irregular blotches
multiplied over the wash, most darker, some lighter, and a dozen dark specks for the
granulation; five to eight darker brush strokes running out from the base; the outline
stroked $1.6\,\ell\,$bp wide for a petal $\ell\,$cm long and multiplied, of which only the inner half shows, for the pigment
that pools at the edge of a wash; nine to thirteen white streaks lifted out in a fan from
the base; and a few fine dark veins. There is no outline in the usual sense.
\end{construction}

\begin{construction}[Tapered strokes]
Every streak, vein and rib is a filled shape rather than a stroke, so that it tapers. For a
cubic B\'ezier with control points $P_0,\dots,P_3$ and widths $w_0$, $w_1$ at its ends, push
each $P_i$ along its normal by $\pm\frac12 w(i/3)$, with the normals at $P_1$ and $P_2$ taken
from the chords $P_2-P_0$ and $P_3-P_1$. The two offset curves joined at the ends are a
tapered stroke in two B\'ezier segments, five times smaller in the file than a sampled
ribbon and indistinguishable from one for a gentle curve.
\end{construction}

\begin{construction}[Flowers]
A \emph{flower} is five petals at $72^\circ$ intervals with $\pm6^\circ$ of jitter and lengths
in $[0.9,1.06]$, then a heart of ten white rays, a white disc, fourteen stamens and a copper
eye. A quarter of the flowers are painted a shade deeper. A \emph{lily} is six pointed
petals, $w(v)=2W\sin^{0.85}(\pi v^{0.8})(1-0.3v)$, arranged in a star and then squashed
vertically by a factor in $[0.45,0.62]$: the star is seen at a tilt, so the petals that point
away are short, and since they are drawn first and darker they read as the back of the
flower. A \emph{blossom} is the same idea at a few millimeters, and a \emph{floret}, the
smallest, is five petals filled as one shape around a dark eye; below $0.7\,$mm in radius it is a
ten-pointed star. Buds, bellflowers and sepals are blades on the same spine.
\end{construction}

\begin{construction}[Leaves]
A leaf of length $1$ has the envelope $e(v)=2.3\,W v^{0.55}(1-v)^{0.95}$. An entire leaf is
that envelope with fine teeth. A lobed leaf, like acanthus, is a narrow body of half-width
$(1-0.75\,d)\,e(v)$, where $d$ is the depth of the cut, together with three to five pointed
lobes on each side; each lobe is itself a blade that leaves the midrib at
$38^\circ$--$58^\circ$ toward the tip and reaches the envelope. Body and lobes are filled as
one shape, under the nonzero rule, so there is no seam. One half of the leaf is multiplied
darker, as if it were folded along its rib; then come the blotches, the pooled edge, and a
white midrib with a white vein along every lobe.
\end{construction}

\begin{construction}[Sprays]
A \emph{cyme}, the inflorescence of the forget-me-not, is a stem whose heading turns by
$\pm\tau(s/L)^{2.3}$ at arc length $s$, with $\tau\in[170^\circ,270^\circ]$: straight at first,
it curls like a scorpion's tail at the end. Eleven to seventeen florets crowd along it,
shrinking into buds toward the curl, and a second, shorter curl forks off it three times in
five. A \emph{sprig} is a stem that forks two or three ways four times, each fork $0.58$ to
$0.74$ the length of its parent, with florets and buds at the tips. An \emph{umbel} fans nine
to thirteen rays over $110^\circ$, each ending in a small umbel of starry florets.
\end{construction}

\subsubsection{Composition}

The garden is a plant, not an arrangement. It grows from one root in the bottom-left
corner of the page, where the base of the stems curls into a scroll. Three stems run up the
column through hand-picked control points, interpolated by Catmull--Rom splines and jittered
by a millimeter or so on each page: the tallest ends in the large flower at the top, one in
a lily, one in a flower half-way up. A fourth stem leaves the root along the bottom of the
page, carrying small leaves that thin out toward its end, and runs out into a copper
hairline; it is the footer. The header is a printer's ivy leaf in copper whose stem becomes
the head rule.

Leaves hang at chosen fractions of the stems' length, alternately left and right. Every
other motif sits at a chosen point and is joined to a chosen stem by a \emph{branch}: a cubic
B\'ezier that leaves the stem along the stem's own direction turned $28^\circ$ to $48^\circ$
toward the target, and arrives along the motif's axis. Because a branch starts tangent to
its parent and ends tangent to its flower, every flower can be traced back to the root, and
that, more than any detail of the painting, is what makes the margin look grown. Behind
everything are one or two \emph{ghosts}: flowers painted with every tone lifted toward the
paper, $t\mapsto t+0.62(1-t)$, and without their fine detail, which give the column depth.

\begin{remark}[Procedural is not random]
The generator is seeded from the page number before the garden is drawn. Page $n$
therefore receives the same garden on every compilation, and a different one from page
$n+1$. The variation the reader notices between pages is real; the reproducibility is what
makes the document stable enough to typeset.
\end{remark}

\begin{proposition}
Let $P$ be the number of primitive instances on a page and $S$ the mean number of path
points per instance. The cost of a page is $\Theta(PS)$ floating-point operations in Lua and
$\Theta(PS)$ bytes of output, and \TeX\ does no arithmetic at all; in practice a page takes
well under a second, for $P\approx 400$ and $S\approx 100$.
\end{proposition}

\begin{proof}[Proof sketch]
Every coordinate is computed once in Lua, transformed to the page, and printed with one
decimal of a PDF point into a single literal node, which \TeX\ only carries to the page.
The same curves can be computed in \TeX\ itself with \texttt{pgfmath}, but it does
arithmetic by macro expansion in \TeX\ dimensions, and a page then takes about a minute. The
size of the output is the real constraint: about
three megabytes of path operators per page, which compress to under one.
\end{proof}

\begin{exercise}
The composition above is authored. Replace it by a growth process in the spirit of Wong,
Zongker and Salesin (\emph{Computer-generated floral ornament}, SIGGRAPH 1998): starting
from the stems, repeatedly find the largest free disc in the margin that is within reach of
a stem, choose a primitive that fits it, and grow a branch to it, until no disc larger than
a threshold remains. The ornament then adapts to any text block, and a page whose text
ends early grows a garden into the space below it.
\end{exercise}

%% text-growth.tex -- the growth stage (used by florilegium-grow.tex only)
\subsubsection{The growth stage}

This document is the answer to the exercise. Its margin is not the authored composition of
the previous section: nothing in it was placed by hand. A Lua program decides where every
stem, leaf and flower goes and how large it is, and then calls the same primitives to paint
them.

\begin{construction}[Stems]
The root is placed near the bottom-left corner of the page. Three stems rise from it; the
$x$ coordinate of each is a smooth function of the height,
\[
  x(y)=c+A\,n\bigl(y/\lambda+\varphi\bigr),
  \qquad A=0.2\,\text{(column width)},\ \lambda\in[6,9]\,\text{cm},
\]
with $n$ a two-octave gradient noise and $\varphi$ a phase drawn per stem, blended into the
root over the first four centimeters. Different phases make the stems sway apart and cross
now and then, as real stems do. The first reaches the top of the page and ends in the
largest flower; the second stops at $60$ to $72\%$ of the height and ends in a lily; the
third at $36$ to $48\%$, in a flower. A footer stem leaves the root along the bottom of the
page. Leaves are hung at the nodes, every $1.4$ to $2.2\,$cm, alternately left and right,
$35^\circ$ to $60^\circ$ off the stem, $2.4\,$cm long at the bottom of the page and $1.3\,$cm at
the top, and only where there is room.
\end{construction}

\begin{definition}
A \emph{proxy} for a motif is a disc of radius $r(s)$, where $s$ is the motif's scale, that
approximates the area it occupies. A point $p$ of the free region $R$ has \emph{free
radius}
\[
  f(p)=\min\Bigl(\,d(p,\partial R),\ \min_i\bigl(\lVert p-c_i\rVert-\kappa\,r_i\bigr)\Bigr),
\]
the radius of the largest disc about $p$ that stays inside $R$ and overlaps no proxy
$(c_i,r_i)$ by more than the fraction $1-\kappa$. With $\kappa=0.85$, a leaf may tuck under
the edge of a flower; stems are proxies too, as chains of small discs.
\end{definition}

\begin{construction}[Growth]
The region $R$ is the margin column, a band along the bottom of the page, and on the last
page also the rectangle between the end of the text and the footer. On a grid of candidates
spaced $2\,$mm apart, each point keeps $f(p)$ and its distance to the nearest stem. Then,
until the largest free radius falls below $3.6\,$mm:
\begin{enumerate}
  \item take the point $p$ that maximizes $f(p)$ among those within $3\,$cm of a stem, with
        a random factor in $[0.85,1.15]$ so that ties break differently on every seed;
  \item choose a motif whose smallest proxy fits, by weight; a motif that has already
        appeared $k$ times has its weight divided by $1+k/c$, and only flowers, lilies,
        umbels and bells have a hard cap. Set its scale so that $r(s)\approx0.95\,f(p)$;
  \item give it an axis pointing away from the nearest stem point, leaning upward for the
        kinds that reach for the light, and grow a branch to it from that point, as in the
        authored composition. The branch becomes a stem for later branches.
\end{enumerate}
Because the largest hole is always taken first, flowers and lilies are placed early,
sprays and leaves fill the middle, and a last pass puts pairs of leaflets, buds and small
clusters of florets on the stems wherever a hole of more than $1.7\,$mm is left next to
one; a hole far from any stem may get one of six copper accents, and otherwise stays empty.
That ordering is the whole reason the result looks \emph{composed} rather than scattered.
\end{construction}

\begin{remark}[Adapting to the text]
The free region is computed from the page layout, so changing the margins moves the
garden. The end of the text is measured where it happens: \verb|\florilegiumtextend|
records \verb|\pagetotal| and \verb|\pagedepth|, the height of the material already on the
current page, and the background of that page is drawn after it. On the last page, the
garden therefore spills into the space below the words in a single Lua\LaTeX\ run: branches
arch out of the main stems into that space and each ends in a flower.
\end{remark}

\begin{remark}[Visitors, marks and seasons]
Butterflies and bees take holes like any motif but need no branch; a snail is put on a
main stem and, on about one page in five, a bird on a twig of its own. Their decisions come
from a random stream of their own, so a page without them is the page it would have been.
A mark in the text, left by \verb|\florilegiumbloom| or by a section heading, is found in
the page box at shipout and gets its flower before any leaf is hung. An age between $0$
and $1$ reweights the motifs from buds in spring to seed heads in autumn.
\end{remark}

\begin{remark}[Cost]
The layout is cheap: the grid has a few thousand points, a proxy only updates the points
near it, and the whole loop takes a fraction of a second. This document of five pages
compiles in about five seconds.
\end{remark}

%% text-sharing.tex -- shared closing section
\subsubsection{Building and sharing}

The package is \texttt{florilegium.sty} with six Lua files beside it: \texttt{ink.lua}
(geometry, color and PDF output), \texttt{flora.lua} (the primitives),
\texttt{authored.lua} and \texttt{grow.lua} (the two compositions), \texttt{initial.lua}
(the decorated initials) and \texttt{florilegium.lua} (the glue to \TeX). It needs Lua\LaTeX\ and one compilation, at
about a second per page. To share it, the three routes that work are, in increasing order
of effort: paste the files into a blank Overleaf project with the compiler set to
Lua\LaTeX\ and share the link; put them in a Git repository with a continuous integration
job that compiles the PDFs on every push and attaches them to tagged releases; or publish
the package so that any document can \texttt{\textbackslash usepackage} the garden and
choose its own seed, column and hue.

Xe\LaTeX\ cannot run Lua, so \texttt{texlua florilegium-xe.lua} draws outside \TeX\ with the
same files, and \texttt{xdvipdfmx} places the pages it writes. Either \verb|xelatex --shell-escape|
draws each item as it is placed, or \texttt{xelatex}, \texttt{texlua} and \texttt{xelatex} do it in
three steps; marks need a second run. The drawings equal Lua\LaTeX's byte for byte.

\begin{remark}
Two mistakes cost time and are worth recording. A literal written into the page bypasses
TikZ, so it also bypasses TikZ's transformations: a primitive called inside
\verb|\begin{scope}[shift=...,rotate=...]| must read pgf's current matrix itself (it is in
\verb|\pgf@pt@aa| to \verb|\pgf@pt@bb| and the dimensions \verb|\pgf@pt@x|,
\verb|\pgf@pt@y|) and apply it in Lua. And a Lua function that returns two values, placed
in the middle of a table constructor such as \verb|{x0, y0, f(t), x1, y1}|, contributes only
its first value; the rest of the table shifts, and the error surfaces three functions
later as arithmetic on \texttt{nil}.
\end{remark}

\florilegiumtextend
\end{document}
